\documentclass[10pt]{article}

\usepackage[utf8]{inputenc}
\usepackage[T1]{fontenc}
\usepackage{geometry}
\geometry{margin=0.9in}
\usepackage{setspace}
\setstretch{0.95}
\usepackage{parskip}
\usepackage{hyperref}
\hypersetup{colorlinks=true,linkcolor=blue,citecolor=blue,urlcolor=blue}
\usepackage{tabularx}
\usepackage{booktabs}
\usepackage{xcolor}

\pagenumbering{gobble}

\begin{document}

\begin{center}
{\Large\bfseries TrayGuard --- Executive Summary}\\[4pt]
{\small Final Design Review \hfill EC ENGR 180DW --- Systems Design}\\[2pt]
{\small Prepared by: Owen Fan, Matthew Li, Andy Ma, Jacob Moore, Qiyu Xin}\\[2pt]
{\small \today}
\end{center}

\section*{Requirements}

\paragraph{Problem.}
Tray-content errors during surgical instrument reconstruction are a recurring
patient-safety risk in sterile processing departments (SPDs). Every hospital
with reusable surgical instruments depends on manual tray assembly: technicians
inspect, identify, sort, and build procedure-specific sets against tray-specific
count sheets. Published evidence reports roughly 9 defects per 100 surgical
cases, with 55\% originating during the assembly step and 44\% of packaging
errors involving wrong instrument specifications. The resulting OR delays
average 10~minutes per affected case, at an estimated annual cost of
\$6.75--9.42~million at a single hospital.

\paragraph{Stakeholders.}
Patients rely on sterile, functional instruments during surgery. Surgeons and OR
teams need complete sets to maintain procedure flow. SPD technicians perform
high-volume, detail-intensive work under production pressure. SPD supervisors
and educators manage training, quality assurance, staffing, and competency
programs. Hospitals absorb financial losses from delayed cases and instrument
rework.

\paragraph{Root cause analysis.}
Root-cause analysis identified six contributing categories --- instrument
design and nomenclature (Product), insufficient structured training
(Knowledge), single-pass assembly with no independent check (Process),
incomplete count sheets and loaner-tray information gaps (Information),
production pressure and ergonomic constraints (Environment), and decentralized
governance with weak error feedback loops (Policy \& Feedback). Three
structural drivers amplify every category: SPD evolved from materials
management rather than patient care, hospital accounting treats SPD as a cost
center without revenue visibility, and no information architecture connects
downstream OR outcomes to upstream SPD actions.

\paragraph{Quantified value.}
TrayGuard targets the training gap among the six root causes.
The study's acceptance bar is a 20~percentage-point pre/post improvement
in simulated tray-assembly accuracy --- consistent with Ofstead et al.~(2023),
who found a \~{}43~point gain (41\% to 84\%) in a structured SP training pilot
with certified professionals. TrayGuard's self-directed simulation is expected
to produce a smaller but visible effect, making 20~points a reasonable
formative threshold. Meeting this bar would provide a replicable preparation
step before supervised clinical training --- reducing clinical supervisor burden
and improving first-pass assembly quality.

\paragraph{Negative externalities.}
Printed AprilTag cards and study flashcards are consumable physical media;
discarded cards contribute to material waste. Accrediting bodies that rely on
documented supervised practice hours may face questions about whether simulated
scores satisfy competency documentation requirements, creating administrative
burden for hospitals that choose not to adopt the system. Controls include
reusable card materials, open module formats, and explicit labeling of results
as simulated practice evidence rather than competency determinations.

\section*{System Design}

\paragraph{Final system.}
TrayGuard is a local tray training and assessment platform. The core product
combines a web application with a file-backed tray module defining instruments,
aliases, required counts, distractors, lookalike pairs, distinguishing
features, and assessment variants for one local tray.

\paragraph{Learning loop.}
The seven-step learner flow is: intake $\rightarrow$ pre-test $\rightarrow$
retrieval-first study cards $\rightarrow$ identification quiz $\rightarrow$
AprilTag-card simulated tray sort $\rightarrow$ post-test $\rightarrow$ metrics
export. The pre-test and study-card phases use browser-based recall; the
simulated tray sort uses physical AprilTag cards placed on a tabletop, scanned
by the device camera, and scored against the tray template.

\paragraph{Subsystems.}
The content module defines each tray as a structured YAML file. The learning
workflow engine orchestrates the seven-step sequence. Retrieval-practice cards
show instrument images with aliases, quantities, and distinguishing features.
The AprilTag scanner detects marker identities and positions from a camera
frame. The scoring engine classifies each item into correct, missing, extra,
wrong, misidentified (lookalike substitution), or wrong-count categories.
Confidence ratings (1--5 scale and ``not sure'' flag) are captured per item.
The export subsystem produces JSON and CSV outputs with hashed learner
identifiers, error categories, and module version fields.

\paragraph{Costs and benefits.}
Direct material costs are \$5--55 for printed cards and a tray surface. The
primary cost is labor: module authoring, prototype implementation, study
moderation, and analysis. The benefit to stakeholders is a repeatable, low-cost
simulated practice environment that surfaces weak-item evidence for instructors
and reduces the number of clinical supervisor hours needed for basic tray familiarization.

\section*{Project Plan}

\paragraph{Original plan versus actual execution.}
The PDR proposed a computer-vision tray checker that would detect instruments
on a tray and compare results against a digital count sheet. Seven staged
detection experiments (800--80~lumens, matte and reflective backgrounds,
separated and overlay layouts) demonstrated that YOLO11s detection degrades
under every variation SPD workflows present, with occlusion mAP50-95 falling to
0.389. Real-instrument experiments confirmed position-cue dominance over shape
learning.

\paragraph{Pivot.}
At Week~6 the project pivoted to a training-centered system with a stronger
near-term evidence path. The new direction changed the proof obligation from
deployment-grade recognition to measurable learning on a simulated local
tray --- a claim that matches the available resources and produces actionable
evidence for a future CV support layer.

\paragraph{Delivered.}
The current prototype includes: a FastAPI web app implementing the full
seven-step learning loop, eight file-backed tray modules, a card-generation
pipeline with AprilTag markers and double-sided study flashcards, a scoring
engine with six error categories, JSON/CSV export with hashed identifiers, and
a GitLab-organized project plan with epics, issues, and milestones.

\paragraph{Deferred.}
Deployment-grade CV recognition, clinical workflow integration, and SPD-site
validation remain out of scope. These are future work items that depend on
validated local content, site-specific training data, and institutional
approval.

\paragraph{FDR placement.}
This Final Design Review documents the complete design process from the
original CV concept through the training pivot to the current prototype,
including experimental evidence for both the abandoned direction and the
selected path. The report captures the engineering rationale, subsystem design,
analysis, evaluation plan, and bounded claims that position TrayGuard as a
validated competency-maintenance tool within the SPD training ecosystem.

\end{document}
